\chapter{Полный 42-скачковый GKSL и его неподвижная алгебра}
\label{ch:v8-full-noise-42-jump-gksl-fixed-algebra-gate}

На полном самосопряжённом кадре $\{F_a\}_{a=1}^{42}$ определим
\begin{equation}
 \mathcal L_{42}(X)=\sum_{a=1}^{42}
 \left(F_aXF_a-\frac12\{F_a^2,X\}\right).
\end{equation}
Это точный конечномерный GKSL-генератор на
$M_{11}(\mathbb C)\oplus M_{10}(\mathbb C)$. Он сохраняет след, единицу и
концевую алгебру.

Прежние 25 операторов скачка принадлежат линейному линейная оболочка нового кадра. Их
совместный коммутант уже был равен $\mathbb CI_{21}$; добавление ещё
семнадцати самосопряжённых направлений может только уменьшить коммутант,
тогда как скаляры остаются в нём. Поэтому
\begin{equation}
 \operatorname{Fix}(\mathcal L_{42})=\mathbb CI_{21}.
\end{equation}
Полный процесс примитивен. Переход к следово-двойственному ортонормированному кадру
является обратимой заменой базиса на том же 42-мерном линейная оболочка и не меняет эту
неподвижную алгебру.

\begin{theorem}[Полная шумовая марковская сшивка]
Сорок два двойственных полю направления скачков задают унитальную сохраняющую след
GKSL-полугруппу с единственной неподвижной алгеброй $\mathbb CI_{21}$.
Примитивность сохраняется при следово-двойственной ортонормировке полного кадра.
\end{theorem}

\begin{nogocriterion}[Математическая полнота не равна физическому выбору]
Нельзя считать равные исходные коэффициенты или следово-двойственную метрику физически
обязательными без независимого Рис/принципа среды. Гейт закрывает носитель,
полную положительность и неподвижная алгебра, но не абсолютное время и подготовку
среды.
\end{nogocriterion}

\begin{protocol}
\begin{enumerate}
 \item операторов скачка: \textbf{$42$};
 \item прежний подкадр: \textbf{$25$};
 \item добавлено: \textbf{$17$};
 \item форма GKSL: \textbf{точно};
 \item сохранение следа и единицы: \textbf{точно};
 \item инвариантность концевых секторов: \textbf{точно};
 \item неподвижная алгебра: \textbf{$\mathbb CI_{21}$};
 \item примитивность: \textbf{получена};
 \item следово-двойственная смена кадра: \textbf{обратима};
 \item физическая скоростная метрика: \textbf{условна};
 \item следующий гейт: \textbf{42-скачковый гамильтониан повторных взаимодействий и
       минимальная среда}.
\end{enumerate}
\end{protocol}

Машинный сертификат сохранён в
\path{s2t_v8_full_noise_42_jump_gksl_fixed_algebra_gate_results.json}.